%% LyX 2.4.4 created this file.  For more info, see https://www.lyx.org/.
%% Do not edit unless you really know what you are doing.
\documentclass[english]{article}
\usepackage[T1]{fontenc}
\usepackage[latin9]{inputenc}
\usepackage{enumitem}
\usepackage{amsmath}
\usepackage{amsthm}
\usepackage{geometry}
\geometry{verbose,tmargin=1in,bmargin=1in,lmargin=1in,rmargin=1in}
\PassOptionsToPackage{normalem}{ulem}
\usepackage{ulem}

\makeatletter
%%%%%%%%%%%%%%%%%%%%%%%%%%%%%% Textclass specific LaTeX commands.
\numberwithin{equation}{section}
\numberwithin{figure}{section}
\newlength{\lyxlabelwidth}      % auxiliary length 

%%%%%%%%%%%%%%%%%%%%%%%%%%%%%% User specified LaTeX commands.
\usepackage{fancyhdr}
\usepackage{lastpage}
\pagestyle{fancy}
\usepackage{enumitem}
\usepackage{ifthen}
\everymath=\expandafter{\the\everymath\displaystyle}
%\DeclareMathSizes{10}{10}{10}{10}
\usepackage{multicol}

\usepackage{tikz}
\usetikzlibrary{patterns}
\usetikzlibrary{plotmarks}
\usepackage{pgfplots}
\pgfplotsset{compat=newest}

\usepackage{calc}
\usepackage[nomessages]{fp}

\usepackage{datetime}
% Added by lyx2lyx
\setlength{\parskip}{\smallskipamount}
\setlength{\parindent}{0pt}

\makeatother

\usepackage{babel}
\begin{document}
\renewcommand{\author}{Dr.\,James Ashe}

\newcommand{\classNum}{232}

\newcommand{\classSec}{A and B}

\newcommand{\examType}{Quiz 3 7.1, 7.2, 7.3, 7.4 and 8.1}

\renewcommand{\date}{Due: \formatdate{8}{4}{2013}}

%%First page's header%%%%%%%%%%%%%%%%%%%%%%
\fancypagestyle{firststyle} %%header settings for first pagr
{
	\fancyhead[L]{MTH \classNum{}(\classSec{}) \author{}\\
	\textbf{\examType{}} ~\date{}}
	\fancyhead[C]{}
	\fancyhead[R]{\textbf{Name:\hspace{4cm}}}
	\setlength{\headsep}{14pt}
}
%%%%%%%%%%%%%%%%%%%%%%%%%%%%%%%%%%%%%%%%%%

%%Next page's header%%%%%%%%%%%%%%%%%%%%%%%
\setlist{nolistsep}
\lhead{\classNum{}(\classSec{})} 
\chead{\textbf{Name:}} 
\cfoot{\thepage{}~of~\pageref{LastPage}} 
\setlength{\headsep}{5pt} 
%%%%%%%%%%%%%%%%%%%%%%%%%%%%%%%%%%%%%%%%%%%

%%Various settings; see the preamble for other settings%%%%%
%\setenumerate[0]{leftmargin=12pt,itemindent=0pt,labelwidth=0pt}
\setlist{topsep=.5pt}
\renewcommand{\labelenumi}{\textbf{\arabic{enumi}.}}
\renewcommand{\labelenumii}{\textbf{\alph{enumii}.}}
\thispagestyle{firststyle}
%%%%%%%%%%%%%%%%%%%%%%%%%%%%%%%%%%%%%%%%%%

\newcommand{\xmax}{0}
\newcommand{\xmin}{-\xmax}
\newcommand{\xstep}{0}
\newcommand{\ymax}{0}
\newcommand{\ymin}{-\ymax}
\newcommand{\ystep}{0} 

\textbf{Justify your answers and show your work.}

\textbf{(a)}\textbf{\emph{ }}\emph{Find the inverse on the specified
interval express it in the form $y=f^{-1}(x)$. }\textbf{(b)}\emph{
State the range and domain of $f^{-1}(x)$. }\textbf{(c)}\textbf{\emph{
}}\emph{Graph $f(x)$ and $f^{-1}(x)$. {[}2 pts each{]}.}
\begin{enumerate}[resume]
\item $f(x)=x^{2}-6x+9$ for $x\geq3$\vspace{1.5cm}\vfill
\end{enumerate}
\emph{Find the derivative of the inverse of the following function
at the specified point on the graph of the inverse function }\emph{\uline{without
finding}}\emph{ $f^{-1}$. }
\begin{enumerate}[resume]
\item \textbf{$f(x)=3x^{5}-2x^{4}+1$}, at $(1,0)$ (point on $f^{-1})$.\vfill
\end{enumerate}
\emph{Use the information about $f(x)$ to answer the following questions.
$f(3)=3$, $f(-4)=\sqrt{5}$, and $f(x)$ has the domain $(-\infty,\infty)$
and range $[\sqrt{5},\infty)$ {[}2pts each{]}.}
\begin{enumerate}[resume]
\item Find $f^{-1}(3)$\vspace{1cm}
\item Find $f^{-1}(\sqrt{5})$\vspace{1cm}
\item What is the domain of $f^{-1}(x)$\vspace{1cm}
\end{enumerate}
\newpage\emph{Perform the indicated operation.}
\begin{enumerate}[resume]
\item $\frac{d}{dx}$$x^{2}e^{-x}$\vfill
\item $\frac{d}{dx}$$\frac{e^{x}+3x}{2}$\vfill
\item $\frac{d}{dx}\left(2\ln x+\sqrt{2}\right)$\vfill
\item $\frac{d}{dx}\left(\ln(3x^{2})-\frac{x}{5}\right)$\vfill
\end{enumerate}
\newpage\emph{Perform the indicated operation.}
\begin{enumerate}[resume]
\item ${\displaystyle \int\left(3e^{-3x}+\frac{1}{x}\right)\,dx}$\vfill
\item $\frac{d}{dx}5^{x}\,dx$\vfill
\item $\frac{d}{dx}x^{\sqrt{8}}$\vspace{1cm}
\item $\frac{d}{dx}x^{\ln x}$ (use logarithmic differentiation)\vfill
\end{enumerate}
\newpage\emph{Perform the indicated operation.}
\begin{enumerate}[resume]
\item ${\displaystyle \int_{1}^{2}2^{x}}$ \vfill
\item ${\displaystyle \int}3xe^{2x}\,dx$\vfill
\item ${\displaystyle \int x\ln2x\,dx}$\vfill
\item ${\displaystyle \int x^{2}\cos4x\,dx}$\vfill
\end{enumerate}

\end{document}
